\documentclass[10pt]{article}
\usepackage[utf8]{inputenc}
\usepackage[T1]{fontenc}
\usepackage{amsmath}
\usepackage{amsfonts}
\usepackage{amssymb}
\usepackage[version=4]{mhchem}
\usepackage{stmaryrd}
\usepackage{bbold}

%New command to display footnote whose markers will always be hidden
\let\svthefootnote\thefootnote
\newcommand\blfootnotetext[1]{%
  \let\thefootnote\relax\footnote{#1}%
  \addtocounter{footnote}{-1}%
  \let\thefootnote\svthefootnote%
}

%Overriding the \footnotetext command to hide the marker if its value is `0`
\let\svfootnotetext\footnotetext
\renewcommand\footnotetext[2][?]{%
  \if\relax#1\relax%
    \ifnum\value{footnote}=0\blfootnotetext{#2}\else\svfootnotetext{#2}\fi%
  \else%
    \if?#1\ifnum\value{footnote}=0\blfootnotetext{#2}\else\svfootnotetext{#2}\fi%
    \else\svfootnotetext[#1]{#2}\fi%
  \fi
}

\begin{document}
\section*{PROBABILITY GU 4155: Spring 2026}
\section*{HOMEWORK \#6 \\
 DO PROBLEMS 1, 3, 4, 9, 12 \\
 TO BE SUBMITTED BY 9:00 PM ON SATURDAY MARCH 14, 2026}
Professor Ioannis Karatzas\\
Department of Mathematics, Columbia University\\
February 3, 2026

Exercise \#1: Lyapunov Inequality. If $X$ is a measurable function on a probability space $(\Omega, \mathcal{F}, \mathbb{P})$ and $0<p<q<\infty$, then

$$
\|X\|_{p} \leq\|X\|_{q}
$$

In particular, for $0<p<q<\infty$ we have $\mathbb{L}^{q} \subseteq \mathbb{L}^{p}$ on a probability space.\\
(This kind of inclusion can fail, and rather dramatically, on a general measure space.)\\
Exercise \#2: Bernštein polynomials. For a given continuous function $f:[0,1] \rightarrow \mathbb{R}$, recall from class the "Bernštein polynomials"


\begin{equation*}
B_{n}(x)=\sum_{k=0}^{n} f(k / n) \cdot\binom{n}{k} x^{k}(1-x)^{n-k}, \quad 0 \leq x \leq 1, \quad n \in \mathbb{N} \tag{0.1}
\end{equation*}


which, as we mentioned in class, converge to $f(\cdot)$ uniformly over the unit interval. Suppose now that the function $f(\cdot)$ satisfies for some $K \in(0, \infty)$ the so-called Lipschitz condition

$$
|f(x)-f(y)| \leq K|x-y|, \quad \forall x, y \in[0,1]
$$

Show that the BernŠtein polynomials of (0.1) satisfy then

$$
\sup _{0 \leq x \leq 1}\left|B_{n}(x)-f(x)\right| \leq \frac{K}{2 \sqrt{n}}, \text { for all } n \in \mathbb{N}
$$

In other words: the Lipschitz condition allows you to get some idea about how fast this uniform convergence happens.

Exercise \#3: An Extended Monotone Convergence Theorem. Suppose $X_{1} \leq \cdots \leq X_{n} \leq X_{n+1} \leq \cdots$ are integrable random variables, monotonically increasing pointwise to the random variable $X=\lim _{n \rightarrow \infty} X_{n}$.

If $\sup _{n \in \mathbb{N}} \mathbb{E}\left(X_{n}\right)<\infty$, show that $X$ is integrable and that we have

$$
\lim _{n \rightarrow \infty} \mathbb{E}\left(X_{n}\right)=\mathbb{E}(X)
$$

Exercise \#4: Riemann-Zeta Calculus. For $s>1$ fixed, let $X: \Omega \rightarrow \mathbb{N}$ be a random variable with the "zeta distribution"

$$
\mathbb{P}(X=n)=\frac{1}{\zeta(s) \cdot n^{s}}, \quad n \in \mathbb{N}, \quad \text { where } \quad \zeta(s):=\sum_{n \in \mathbb{N}} \frac{1}{n^{s}}
$$

is the Riemann zeta-function. For each $m \in \mathbb{N}$, consider the event\\
$A_{m}:=\{\omega \in \Omega: m$ is a factor of $X(\omega)\}=\{\omega \in \Omega: X(\omega)=m k(\omega)$, for some $k(\omega) \in \mathbb{N}\}$.\\
(i) Show that $\mathbb{P}\left(A_{m}\right)=1 / m^{s}$.\\
(ii) With $\mathfrak{P}$ denoting the set of prime numbers, conclude from (i) that the events in the collection $\left\{A_{p}\right\}_{p \in \mathfrak{P}}$ are independent,\\
(iii) Use (i) and (ii) to derive the celebrated Euler formula

$$
\frac{1}{\zeta(s)}=\prod_{p \in \mathfrak{P}}\left(1-\frac{1}{p^{s}}\right)
$$

Exercise \#5: Conditional Independence: (i) Show that if the events $A$ and $B$ are independent, so are $A$ and $B^{c}$; that so are $B$ and $A^{c}$; and that so are $A^{c}$ and $B^{c}$.\\
(ii) Two events $A, B$ are said to be conditionally independent given a third event $C$ with $\mathbb{P}(C)>0$, if

$$
\mathbb{P}(A \cap B \mid C)=\mathbb{P}(A \mid C) \cdot \mathbb{P}(B \mid C)
$$

When $\mathbb{P}(C \cap B)>0$, argue that this property is equivalent to

$$
\mathbb{P}(A \mid C \cap B)=\mathbb{P}(A \mid C)
$$

Exercise \#6: For a sequence of subsets $\left\{E_{n}\right\}_{n \in \mathbb{N}}$ of the real line, determine the sets

$$
\left\{E_{n}, \text { ev. }\right\}:=\bigcup_{k \in \mathbb{N}} \bigcap_{n \geq k} E_{n}, \quad\left\{E_{n}, \text { i.o. }\right\}:=\bigcap_{k \in \mathbb{N}} \bigcup_{n \geq k} E_{n}
$$

when the sets $E_{n}$ are specified as follows:\\
(i): $E_{n}=\left[a_{n}, b_{n}\right)$ with $0<a_{n}<b_{n}, b_{n}>1$ and $\lim _{n \rightarrow \infty} a_{n}=0, \lim _{n \rightarrow \infty} b_{n}=1$.\\
(ii): $E_{n}=\{m / n: m \in \mathbb{N}\}$ for each $n \in \mathbb{N}$.

Exercise \#7: Suppose the real-valued random variables $X, Y$ are identically distributed and we have $X \geq Y$ pointwise. Show then $X=Y$ a.e.

Exercise \#8: Suppose that $X_{1}, X_{2}, \cdots$ are independent random variables with common, standard normal (Gaussian) distribution. Find a sequence $\left\{\mathfrak{a}_{n}\right\}_{n \geq 2}$ of positive real numbers with $\lim _{n \rightarrow \infty} \mathfrak{a}_{n}=\infty$ and such that

$$
\limsup _{n \rightarrow \infty}\left(\frac{\left|X_{n}\right|}{\mathfrak{a}_{n}}\right)=1 \quad \text { holds almost surely. }
$$

Hint: The "tail" $\mathbb{P}\left(X_{1}>x\right)=\int_{x}^{\infty} \varphi(u) d u, x \in(0, \infty)$ of the standard normal distribution satisfies the so-called Mill's ratio asymptotics

$$
\lim _{x \rightarrow \infty} \frac{\mathbb{P}\left(X_{1}>x\right)}{\varphi(x) / x}=1, \quad \text { where } \quad \varphi(x):=(2 \pi)^{-1 / 2} e^{-x^{2} / 2}
$$

Problem \#9: On a suitable probability space, construct\\
(i) a sequence of random variables that converges in probability, but not almost surely;\\
(ii) a sequence of random variables that converges almost surely, but does not exhibit fast convergence in probability; and\\
(iii) a sequence of random variables that converges almost surely, but not in $\mathbb{L}^{1}$.\\
(Hint: All these sequences can be constructed on the space $\Omega=(0,1]$ with Lebesgue measure $\mathbb{P}$ on its Borel sets.)

Exercise \#10: A strengthening of Borel-Cantelli: If $A_{1}, A_{2}, \cdots$ are pairwise-independent events with $\sum_{n \in \mathbb{N}} \mathbb{P}\left(A_{n}\right)=\infty$, then for

$$
Q_{n}:=\frac{\sum_{j=1}^{n} \mathbf{1}_{A_{j}}}{\sum_{j=1}^{n} \mathbb{P}\left(A_{j}\right)} \quad \text { we have } \lim _{n \rightarrow \infty} Q_{n}=1 \quad \text { a.e. }
$$

(Hint: Establish convergence in probability first. Then proceed to show almost-everywhere convergence: first for the subsequence $\left\{Q_{n_{k}}\right\}_{k \in \mathbb{N}}$ defined via

$$
n_{k}=\inf \left\{n \geq 1: \sum_{j=1}^{n} \mathbb{P}\left(A_{j}\right) \geq k^{2}\right\}, \quad k \in \mathbb{N}
$$

and then for the entire sequence.)\\
Exercise \#11: In his proof of the the Strong Law of Large Numbers, ${ }^{1}$ why does Etemadi use a "pliant" (expanding) Procrustean bed, whose length increases by one unit on each successive day, and not a Procrustean bed of fixed length?

To wit: why does he define, for each $j \in \mathbb{N}$, the random variable


\begin{equation*}
Z_{j}:=h_{j}\left(X_{j}\right), \quad \text { with } \quad h_{j}(u):=u \cdot \mathbf{1}_{[0, j]}(u), \quad u \in[0, \infty) \tag{0.2}
\end{equation*}


and not, say,


\begin{equation*}
\widehat{Z}_{j}:=h\left(X_{j}\right), \quad \text { with } \quad h(u):=u \cdot \mathbf{1}_{[0, M]}(u), \quad u \in[0, \infty) \tag{0.3}
\end{equation*}


for some fixed $M \in \mathbb{N}$ ? What is achieved by the choise (0.2) that fails for the choice (0.3)?

\footnotetext{${ }^{1}$ Discussed in detail in class.
}Exercise \#12: The Concentration of Measure Phenomenon. (i) Suppose $X_{1}, X_{2}, \cdots$ are square-integrable random variables with

$$
m=\mathbb{E}\left(X_{j}\right), \quad \mathbb{E}\left[\left(X_{i}-m\right)\left(X_{j}-m\right)\right]=0, \quad \operatorname{Var}\left(X_{j}\right):=\mathbb{E}\left[\left(X_{j}-m\right)^{2}\right] \leq C
$$

valid for all natural numbers $j$ and $i \neq j$, and some real constants $m$ and $C>0$.\\
(i) Show that we have, both in $\mathbb{L}^{2}$ and in probability:


\begin{equation*}
\lim _{n \rightarrow \infty} \frac{1}{n} \sum_{j=1}^{n} X_{j}=m \tag{0.4}
\end{equation*}


(ii) Suppose now that $Y_{1}, Y_{2}, \cdots$ are independent random variables, and have a common distribution which is uniform (i.e., normalized Lebesgue measure) on the interval ( $-1,1$ ). Argue that


\begin{equation*}
\lim _{n \rightarrow \infty} \frac{1}{n} \sum_{j=1}^{n} Y_{j}^{2}=\frac{1}{3} \tag{0.5}
\end{equation*}


holds in probability. (Does this hold also in a stronger sense, perhaps?)\\
(iii) Consider now the cube $\mathbb{K}^{n}:=(-1,1)^{n}$ in Euclidean space $\mathbb{R}^{n}$. Justify the statement:\\
"For n large, most of the volume of this cube is concentrated very near the sphere (i.e., the boundary of the ball) of radius $\sqrt{n / 3}$ in $\mathbb{R}^{n}$ ".\\
Do this, by establishing for every $\varepsilon \in(0,1)$ the property


\begin{equation*}
\lim _{n \rightarrow \infty} \frac{1}{2^{n}} \cdot \lambda_{n}\left(A_{n, \varepsilon} \cap \mathbb{K}^{n}\right)=1 \tag{0.6}
\end{equation*}


Here $\lambda_{n}$ is Lebesgue measure (i.e., Euclidean volume) in $\mathbb{R}^{n}$, so $\lambda_{n}\left(\mathbb{K}^{n}\right)=2^{n}$; and


\begin{equation*}
A_{n, \varepsilon}:=\left\{y \in \mathbb{R}^{n}: \frac{n}{3}(1-\varepsilon) \leq y_{1}^{2}+\cdots+y_{n}^{2} \leq \frac{n}{3}(1+\varepsilon)\right\} \tag{0.7}
\end{equation*}


is a "thin crust" around the sphere of radius $\sqrt{n / 3} \cdot{ }^{2}$

\footnotetext{${ }^{2}$ In dimension $n=300$ and with $\varepsilon=.02$, the thin crust becomes $A_{n, \varepsilon}=\left\{y \in \mathbb{R}^{n}: 9.9 \leq\|y\| \leq 10.1\right\}$, so the part $\mathbb{K}^{n} \backslash A_{n, \varepsilon}$ of the cube outside the thin crust is essentially "empty"!
}
\end{document}
